% ===== PRÉAMBULE CANONIQUE — à coller VERBATIM en tête de CHAQUE .tex =====
\documentclass[11pt,a4paper]{article}
\usepackage[utf8]{inputenc}\usepackage[T1]{fontenc}\usepackage{lmodern}
\usepackage[french]{babel}
\usepackage{amsmath,amssymb,mathtools}\usepackage{icomma}
\usepackage[version=4]{mhchem}          % \ce{...} : équations & formules
\usepackage{siunitx}\sisetup{locale=FR} % \qty{}{}, \si{}, \ang{}
\usepackage{chemfig}                    % \chemfig{...} : structures développées
\usepackage{xcolor}\usepackage{colortbl}
\usepackage{tikz}\usepackage{pgfplots}\pgfplotsset{compat=1.18}
\usepackage[most]{tcolorbox}\usepackage{enumitem}\usepackage{array,booktabs}
\usepackage[a4paper,margin=2.2cm]{geometry}
\usetikzlibrary{arrows.meta,calc,angles,quotes,positioning,patterns,backgrounds,shapes.geometric}
\definecolor{primary}{RGB}{222,49,99}\definecolor{primarydark}{RGB}{150,22,55}
\definecolor{defcol}{RGB}{0,114,178}\definecolor{methcol}{RGB}{52,92,148}
\definecolor{excol}{RGB}{0,135,90}\definecolor{attncol}{RGB}{200,40,55}
\tcbset{boxbase/.style={enhanced,breakable,boxrule=0.9pt,arc=2.5pt,
  left=3mm,right=3mm,top=2.3mm,bottom=2.3mm,fonttitle=\bfseries,coltitle=white}}
% --- boîtes TUTORIEL (noms EXACTS, ne pas renommer) ---
\newtcolorbox{cle}[1][]{boxbase,colback=defcol!6,colframe=defcol,
  title={\textsc{À retenir}\ifx\relax#1\relax\relax\else\textmd{ : \itshape #1}\fi}}
\newtcolorbox{methode}[1][]{boxbase,colback=methcol!7,colframe=methcol,sharp corners,
  title={\textsc{Méthode}\ifx\relax#1\relax\relax\else\textmd{ --- \itshape #1}\fi}}
\newtcolorbox{exemplebox}[1][]{boxbase,colback=excol!5,colframe=excol,
  title={\textsc{Exemple}\ifx\relax#1\relax\relax\else\textmd{ : \itshape #1}\fi}}
\newtcolorbox{piege}[1][]{boxbase,colback=attncol!6,colframe=attncol,
  title={\textsc{Piège}\ifx\relax#1\relax\relax\else\textmd{ : \itshape #1}\fi}}
\newtcolorbox{remarque}[1][]{boxbase,colback=black!4,colframe=black!45,
  title={\textsc{Remarque}\ifx\relax#1\relax\relax\else\textmd{ : \itshape #1}\fi}}
% --- boîtes EXERCICE / CORRIGÉ (noms EXACTS, auto-comptées) ---
\newtcolorbox[auto counter]{exo}[1][]{boxbase,colback=primary!4,colframe=primary,
  title={\textsc{Exercice}~\thetcbcounter\ifx\relax#1\relax\relax\else\textmd{ --- \itshape #1}\fi}}
\newtcolorbox[auto counter]{corrige}[1][]{boxbase,colback=excol!4,colframe=excol,
  title={\textsc{Corrigé}~\thetcbcounter\ifx\relax#1\relax\relax\else\textmd{ --- \itshape #1}\fi}}
\newcommand{\R}{\mathbb{R}}\newcommand{\N}{\mathbb{N}}\newcommand{\Z}{\mathbb{Z}}\newcommand{\ds}{\displaystyle}
% (un \usepackage{fancyhdr}+en-tête est possible ; il est ignoré côté web)
% =========================================================================
\begin{document}
\newcommand{\AX}[2]{\ensuremath{\mathrm{AX_{#1}E_{#2}}}}
\begin{corrige}[écrire la notation \AX{n}{m}]
\begin{enumerate}[label=\alph*)]
  \item \ce{CH4} : \AX{4}{0}.
  \item \ce{NH3} : \AX{3}{1}.
  \item \ce{H2O} : \AX{2}{2}.
  \item \ce{CO2} : \AX{2}{0}.
  \item \ce{BF3} : \AX{3}{0}.
\end{enumerate}
\end{corrige}

\begin{corrige}[nombre de sites $\to$ géométrie de répulsion]
\begin{enumerate}[label=\alph*)]
  \item $2$ : \textbf{linéaire}.
  \item $3$ : \textbf{triangulaire} (trigonale plane).
  \item $4$ : \textbf{tétraédrique}.
  \item $5$ : \textbf{bipyramide trigonale}.
  \item $6$ : \textbf{octaédrique}.
\end{enumerate}
\end{corrige}

\begin{corrige}[\AX{n}{m} $\to$ géométrie moléculaire]
\begin{enumerate}[label=\alph*)]
  \item \AX{2}{0} : \textbf{linéaire}.
  \item \AX{3}{0} : \textbf{trigonale plane}.
  \item \AX{4}{0} : \textbf{tétraédrique}.
  \item \AX{3}{1} : \textbf{pyramidale} (à base triangulaire).
  \item \AX{2}{2} : \textbf{coudée}.
\end{enumerate}
\end{corrige}

\begin{corrige}[compter les sites de répulsion]
Une liaison multiple compte pour $1$ ; un doublet libre compte pour $1$.
\begin{enumerate}[label=\alph*)]
  \item \ce{CO2} : $2$ doubles $= \mathbf{2}$ sites (\AX{2}{0}).
  \item \ce{SO2} : $2$ liés $+ 1$ doublet $= \mathbf{3}$ sites (\AX{2}{1}).
  \item \ce{HCN} : $1$ triple $+ 1$ simple $= \mathbf{2}$ sites (\AX{2}{0}).
  \item \ce{BF3} : $3$ liaisons $= \mathbf{3}$ sites (\AX{3}{0}).
  \item \ce{H2O} : $2$ liaisons $+ 2$ doublets $= \mathbf{4}$ sites (\AX{2}{2}).
\end{enumerate}
\end{corrige}

\begin{corrige}[l'angle idéal]
\begin{enumerate}[label=\alph*)]
  \item linéaire : $\mathbf{180^\circ}$.
  \item trigonale plane : $\mathbf{120^\circ}$.
  \item tétraédrique : $\mathbf{109{,}5^\circ}$.
  \item octaédrique : $\mathbf{90^\circ}$.
\end{enumerate}
\end{corrige}

\end{document}
